% ===== 原文 pp.98--102（tr-07c）=====
% 接 tr07b（p.97 页末，句子在页界处中断）；本片段首句完成该跨页句。

% [原文 p.98]
$T(\fg,C)$，其上带有一个由沿一个方向改变 4d 规范耦合所定义的 Janus 畴壁。这一方向上的进一步工作包括 [698, 711]。

正如自 Witten 的 [738] 以来人们所知，许多纽结不变量可以表示为规范理论的配分函数或其他可观测量。参考文献的样本见 [692, 739--746] 与综述 [747]，以及通过研究超多项式族来计算纽结不变量的一系列工作 [748--762]。

另一种不同的拓扑扭曲把纽结与三维流形的同调不变量（单极子/Heegaard Floer 同调与 Khovanov--Rozansky 同调）用 3d $\cN=2$ 理论 $T(\fg,C_3,K_1)$ 在黎曼面上的部分拓扑扭曲来表达 [698, 706, 716]。全息计算 [717, 719--722, 763--766] 在大 $N$ 下探测或使用这一对应。从 4d $\cN=2$ 超共形指标到 3d $\cN=2$ 球面配分函数的维数约化，对应着从 2d $q$-YM 到双曲流形的维数氧化（dimensional oxidation）[675, 767]。3d/3d 对应也可以用高形式对称性来加细 [724]。半 BPS 的 2d $\cN=(0,2)$ 边界条件与 3d $\cN=2$ 理论的畴壁在 [609] 及后续论文中作了研究，其一个分支就是下面要讨论的 2d/4d 对应 [112]。
% [Tutor Note: p.98 原文把 oxidation 拼作 "oxydation"（旧式拼法），这里照译并保留英文原词。
% 其含义是维数约化的逆过程：4d 超共形指标在 $S^3$ 上的约化给出 3d 球面配分函数，
% 反过来就是从 2d $q$-YM "氧化"到双曲流形上的理论。]

\subsection{更多的几何设置}

\para{2d/4d 对应。}
把 6d $(2,0)$ 理论在 $T^2\times C_4$ 上约化、沿四维流形 $C_4$ 作部分拓扑扭曲，就得到 2d $\cN=(0,2)$ 超对称规范理论。这一设置研究得较少，因为四维流形的拓扑比 4d/2d 与 3d/3d 对应中的情形复杂得多。M5 膜作用量所需的相关扭曲已在 [768, 769] 中显式构造出来。

[112] 中提出了一个 à la AGT 的字典：$C_4$ 上的 Vafa--Witten 配分函数 [770] 就是 2d $(0,2)$ 理论的超共形指标，而 4d 的 Kirby 演算则对应着 2d $\cN=(0,2)$ 理论之间的对偶，例如 [771]。带边界 $\partial(C_4)$ 的四维流形，对应由 3d/3d 对应与 $\partial(C_4)$ 相关联的 3d $\cN=2$ 理论之间的畴壁 [609]（另见 [772]）。形如 $\C P^1\times C$ 的四维流形在 [773] 中被考虑，并提供了 class S 理论的一个类似物。这可以通过插入 6d $(2,0)$ 理论的缺陷进一步丰富。有若干变体：紧化在 $S^2\times C_4$ 上 [774]、推广到 6d $(1,0)$ 理论 [775]，以及使用不同的扭曲把四维流形不变量与 2d $\cN=(0,2)$ 手征关联函数联系起来 [776, 777]。阿贝尔情形在 [778] 中有详细研究。另见 [779, 780]。

\para{更少的超对称。}
人们可以从 4d $\cN=2$ 理论的破超对称形变以及 S 对偶 [139, 170, 781--787] 中了解强耦合 4d $\cN=1$ 理论的性质，例如 Seiberg 对偶的类似物。这些进展导致为不承认 4d $\cN=2$ 拉氏量描述的 4d $\cN=2$ AD 理论找到了 4d $\cN=1$ 拉氏量描述，例如 [410, 788--801] 所做的那样（另见超对称更多的 [802]）。它们也导致新的 4d $\cN=2$ 理论，这些理论可能在具有最小中心荷 $(a,c)$ 的意义上是"极小"的 [797]。

获得 4d $\cN=1$ 理论的另一条途径是考虑 6d $\cN=(2,0)$ SCFT 的更一般紧化，这相当于把 M5 膜放在一条复曲线上，
% [原文 p.99]
而这条曲线位于 Calabi--Yau 三维复流形之内 [355, 361, 528, 784, 791, 803--809]。SW 几何的推广似乎存在 [810]。进一步的约化给出 3d $\cN=2$ 理论 [811]。

通向更低超对称的一条特别自然的路径来自 M 理论设置的轨形。位于 $\Z_k$ 奇点的 M5 膜的 6d $(1,0)$ 理论，有许多非常有意思的约化，给出称为 class $S_k$ 的 4d $\cN=1$ 理论 [615, 674, 812--824]；M5 膜探测更一般的 ADE 奇点的情形见 [825--827]。原则上这导致 AGT 对应的类似物，但 $\cN=1$ 理论的 $Z_{S^4}$ 配分函数带有一些歧义，且瞬子配分函数尚不知道（不过可以看一个非常有意思的提议 [240]）。

除这些轨形化的 M5 膜之外，还存在由 F 理论构造出的一大群 6d $(1,0)$ 理论 [828, 829]，综述见 [119, 830]。把它们进一步紧化在环面上给出 4d $\cN=2$ 超对称，重现了许多 class S 理论 [814]。这一设置也曾在黎曼面上 [831--834]，或在带通量的环面上 [820, 825--827] 被研究过，以得到 4d $\cN=1$ 理论；带面算子的情形见 [579, 833, 835]。从 6d 到 5d 的约化也很有意思 [836]。

另外两族丰富的 4d $\cN=1$ SCFT 是二分箭图（bipartite quivers）[837, 838]，以及探测环面奇点定向折叠（orientifold）的 D3 膜 [839--843]。

\para{杂项。}
通过微调（并解析延拓）参数，人们得到极小模型的 AGT 对应 [277, 487, 844--848]，以及数学意义上 AGT 的一个"有限"版本 [849--851]。

Class S 的 4d $\cN=2$ 理论还曾在其他几何上作局域化：$S^2\times S^2$ 对应 Liouville 引力 [852]，$S^2\times S^1\times I$ 对应复 Toda CFT [853]，$S^2\times T^2$ 见 [657]。

不同于把 M5 膜约化在乘积几何上，约化 D3 膜给出 2d/2d 对应 [854--857]，而 M2 膜给出 1d/2d 对应 [858, 859]。[715] 中还提出了一个 3d/3d $\cN=1$ 对应。

\section{结论}

% [Tutor Note: p.99 的 §10 要点框在原书中没有深色标题条（与 §9 等处的 keybox 不同），
% 因此这里直接用无标题的 tcolorbox 复刻同样的浅灰底、黑框样式，而不套用 keybox 环境
% （后者必须带标题）。框内三条即原文的三条要点，顺序与加粗标签照抄。]
\begin{tcolorbox}[colback=black!5,colframe=black!55,boxrule=0.6pt,arc=2pt,
  left=6pt,right=6pt,top=4pt,bottom=4pt,fontupper=\small]
\begin{itemize}[leftmargin=1.4em,itemsep=1pt,topsep=1pt]
  \item \textbf{BPS/CFT 对应}：相当一般的 4d $\cN=2$ 理论的 SW 曲线是可积系统的谱曲线。
  \item \textbf{拓扑弦}给出由约化 $(p,q)$ 五膜网、或由 IIB 几何工程得到的理论的 $Z_{\text{inst}}$。
  \item 带对数势的\textbf{矩阵模型}给出拉氏量 4d $\cN=2$ 理论的 $Z_{\text{inst}}$。
\end{itemize}
\end{tcolorbox}

关于 AGT 对应还有更多可说的。最明显的是，我们尚未把这个对应放进 BPS/CFT 对应的更广阔背景中------即 4d $\cN=2$ 规范理论与 SW 几何背后的可积模型之间的对应。我们也略去了与精细化拓扑弦及矩阵模型之间的联系。

\para{可积系统。}
许多 4d $\cN=2$ 理论的 SW 解可以实现在已知可积系统的谱曲线上 [860--863]，例如周期 Toda 自旋链、
% [原文 p.100]
Calogero--Moser、Ruijsenaars、sine-Gordon 等等。作为一个相当一般的例子，对 class S 的 4d $\cN=2$ 理论，它就是 Hitchin 可积系统 [167]。把 4d 理论放在 $\epsilon_2=0$ 的 Omega 背景上（Nekrasov--Shatashvili 极限）对应于把可积系统量子化，而打开 $\epsilon_2$ 则给出进一步的精细化 [864]。人们也已理解如何在上述讨论中纳入面算子。

Nekrasov 主张把上述考量视为一种 BPS/CFT 对应，综述见 [71, 280, 865--867]（另见 [70, 189, 190, 239, 628, 645, 868--870]），它把具有 8 个超荷的超对称规范理论（例如 4d $\cN=2$）与可积模型及 2d CFT 联系起来。由于这适用于 class S 理论之外，人们可以把 AGT 对应仅仅看作它的一个实例------而在这一实例上人们能取得进一步进展。

另一个实例是 Bethe/规范对应，粗略地说它出现在 Nekrasov--Shatashvili 极限中。在这一极限下，Omega 形变的 4d $\cN=2$ 理论约化为一个 2d $\cN=(2,2)$ 理论，后者的性质与量子可积系统的性质相符。例如，2d 理论的扭曲手征环给出量子 Hamilton 量，超对称真空对应 Bethe 态，而 2d 扭曲超势是可积系统的 Yang--Yang 函数。这一极限在 [864, 871--891] 等文献中作了研究。
% [Tutor Note: p.100 的 Nekrasov--Shatashvili 极限核对：$\epsilon_2=0$ 时 Omega 背景只在一个平面内旋转，
% 规范理论的瞬子展开成为关于 $\epsilon_1$ 的 WKB 型级数，这正对应把 SW 可积系统的量子化参数取作
% $\hbar=\epsilon_1$；打开 $\epsilon_2$（精细化）则在此之上再引入第二个形变参数。
% 锚点：$\epsilon_1\to0$ 的预势展开应重现可积系统的经典谱曲线（即 SW 曲线）。]

关于可积模型与 class S 理论的进一步未经整理的参考文献，见 [57, 211, 279, 292, 400, 524, 530, 561, 584, 593, 609, 633, 634, 748, 864, 883, 892--949]。

\para{拓扑弦。}
拓扑弦及其与 AGT 对应的关系在 [950] 中有综述。

对由 IIB 几何工程 [630, 951, 952] 实现、或作为生活在 $(p,q)$ 五膜网上的 5d $\cN=1$ 理论之维数约化而得到的 4d $\cN=2$ 理论，瞬子配分函数可以表示为拓扑弦的配分函数 $Z_{\text{top}}$，综述见 [953]。正如 [954] 为解释 AGT 对应所主张的，$Z_{\text{top}}$ 还可以进一步用带对数势的 Penner 型矩阵模型来表示 [955--961]，而这些矩阵模型与共形块的 Dotsenko--Fateev 积分表示相符。我们很快会回到这些矩阵模型。

拓扑弦配分函数 $Z_{\text{top}}$ 通过拓扑顶点形式计算：非精细化情形 $\epsilon_1=-\epsilon_2$ 在 [962--965] 中发展出来，而一般的 $\epsilon_1,\epsilon_2$ 有两个表述，分别在 [966, 967] 与 [968, 969] 中给出；它们的等价性在 [970] 中通过把精细化拓扑顶点实现为 Ding--Iohara--Miki（DIM）代数的交织算子（intertwiner）而得到解释。进一步的检验与微妙之处另见 [911, 971--974]，世界叶视角见 [975, 976]，保证 $Z_{\text{top}}$ 与所谓"优选方向"无关的对偶性讨论见 [977, 978]，通过引入新拓扑顶点推广到 A 型箭图之外见 [979, 980]。有时 $Z_{\text{top}}$ 也可以改用全纯反常方程 [879, 953, 981--986]、爆破方程（blowup equations）[987] 或量子曲线 [988, 989] 来自举。

[339--342] 中对 $T_N$ 理论配分函数的计算（因而也是 Toda CFT 三点函数的计算）特别有意思。面算子及其与几何跃迁、$qq$-特征标的关系在 [56, 58, 59, 562, 660, 934, 971, 990--992] 中讨论。基于精细化拓扑顶点的其他 AGT 进展

% [原文 p.101]
包括 [229, 846, 988, 993--1006]。

\para{对称性与特殊多项式。}
AGT 背景下对精细化拓扑弦重新燃起的兴趣，导致了大量推广 Jack 与 Macdonald 多项式的特殊多项式族的出现，包括 Macdonald--Kerov 函数、广义 Schur 函数、椭圆广义 Macdonald 多项式等等：近期发展见 [598, 762, 1007--1020]。

这些发展基于推广 W 代数的各种底层对称代数 [131, 330]，它们值得由更有资格的人单独写一篇综述。我对文献的考察表明，主要角色如下。我发现参考文献 [268, 1021] 是有用的起点。
% [Tutor Note: p.101 起的代数清单是本节最易混淆处，建议对照原书 p.101--102 逐条核对记号：
% $\mathcal{E}_{\sigma,\bar\sigma}$（椭圆 Hall 代数）、DIM、$\ddot{\mathcal U}_{q_1,q_2,q_3}(\mathfrak{gl}_1)$（量子环面代数）、
% $\ddot Y_{h_1,h_2,h_3}(\mathfrak{gl}_1)$（仿射 Yangian）、$\mathrm{SH}^c$（球面退化 DAHA/COHA）。
% 原文的双点记号表示稳定极限/球面版本，与单点不同，抄写时勿漏。]

\begin{itemize}[leftmargin=1.4em,itemsep=4pt]
  \item （中心扩张的）\emph{椭圆 Hall 代数} $\mathcal{E}_{\sigma,\bar\sigma}$ 由 [1022] 引入，是由 $u_{m,n}$（$(m,n)\in\Z^2$，$u_{0,0}=0$）以及相互对易的生成元 $\sigma,\bar\sigma$ 生成的一个结合代数。它们服从的交换关系把 $[u_y,u_x]$（对某些 $x,y\in\Z^2$）表示为 $\sigma,\bar\sigma$ 以及位于连接 $x+y$ 与原点的线段上的生成元 $u_z$（$z\in\Z^2$）的多项式。该代数可以看作球面双仿射 Hecke 代数（DAHA）$\ddot{\mathrm{SH}}_n$ 的稳定极限 [1023] $\mathcal{E}_{\sigma,\bar\sigma}=\ddot{\mathrm{SH}}_\infty$。后者也称为 Cherednik 代数，由 [1024] 引入，综述见 [1025]。

  该代数 $\mathcal{E}_{\sigma,\bar\sigma}$ 允许 $\mathrm{SL}(2,\Z)$ 通过自同构作用，后者由它在 $\Z^2$ 上的作用诱导。对任意互素的 $a,b$，由 $u_{ma,mb}$（$m\in\Z$）生成的子代数构成量子群 $U_q(\mathfrak{gl}_1)$ 的一个拷贝，而这些拷贝在 $\mathrm{SL}(2,\Z)$ 作用下互相交换。

  令 $\mathcal{M}_{\mathrm{U}(N)}$ 为 $\C^2$ 上非交换 U($N$) 瞬子的模空间（Gieseker 框架模空间）。它的等变 K-理论允许 $\mathcal{E}_{\sigma,\bar\sigma}$ 的一个作用 [1021, 1026, 1027]。

  \item \emph{DIM 代数}（由 Ding--Iohara [1028]、Miki [1029] 以及其他人 [1022, 1026] 独立发现）是一个依赖于复参数的结合代数，参数满足 $q_1q_2q_3=1$。它由 $\psi_0^{-1}$ 以及 $e_i,f_i,\psi_i$（$i\in\Z$）生成，并满足诸如
  \[
  [e_i,f_j]=\frac{1}{(1-q_1)(1-q_2)(1-q_3)}\Big((\delta_{i+j>0}-\delta_{i+j<0})\psi_{i+j}+\delta_{i+j=0}(\psi_0-\psi_0^{-1})\Big)
  \]
  这样的二次关系。

  \item \emph{量子环面代数} $\ddot{\mathcal U}_{q_1,q_2,q_3}(\mathfrak{gl}_1)$，也称量子连续 $\mathfrak{gl}_\infty$，是 DIM 代数模掉三次 Serre 关系 $[e_0,[e_1,e_{-1}]]=0$ 与 $[f_0,[f_1,f_{-1}]]=0$ 所得的商。它已出现在 [1029] 中，并在 [1030--1032] 中得到探讨。它也是椭圆 Hall 代数 $\mathcal{E}_{\sigma,\bar\sigma}$ 在 $\bar\sigma=1$ 处的一个特殊化 [1032]；因此，它作用在上述等变 K-理论上。

  \item $\mathfrak{gl}_1$ 的\emph{仿射 Yangian} [269, 1033]，记作 $\ddot Y_{h_1,h_2,h_3}(\mathfrak{gl}_1)$（$h_1+h_2+h_3=0$），由 $e_i,f_i,\psi_i$（$i\ge 0$）生成，其关系由 $\ddot{\mathcal U}_{q_1,q_2,q_3}(\mathfrak{gl}_1)$ 的关系导出。它是球面退化 DAHA 的一个稳定极限。$\mathcal{M}_{\mathrm{U}(N)}$ 的等变上同调允许 $\ddot Y_{q_1,q_2,q_3}(\mathfrak{gl}_1)$ 的一个作用 [268, 1021]。
% [Tutor Note: p.101 末行原文把作用在等变上同调上的代数写成 $\ddot Y_{q_1,q_2,q_3}(\mathfrak{gl}_1)$，
% 而该条开头定义的仿射 Yangian 用的是参数 $h_1,h_2,h_3$（$\ddot Y_{h_1,h_2,h_3}$）。
% 这是原文记号不统一（$q$ 与 $h$ 两种参数化），照译不改；阅读时按"同一个代数的两种参数化"理解即可。
% 依据：同段说明它由 $\ddot{\mathcal U}_{q_1,q_2,q_3}$ 的关系导出，说明此处的 $q$ 参数化指同一类对象。]

% [原文 p.102]
  \item 记作 $\mathrm{SH}^c$ 的代数 [268, 269] 同构于 $\ddot Y_{h_1,h_2,h_3}(\mathfrak{gl}_1)$ 在 $h_1=1$ 处的特殊化 [1033]。它有多种构造，例如 $\mathrm{GL}(\infty)$ 的球面退化 DAHA，或带一个顶点与一个圈的箭图的球面上同调 Hall 代数（COHA）。
  在拓扑完备化之下，$\mathrm{SH}^c$ 在参数的特定取值处与 $W_N$ 手征代数的泛包络代数相符 [268]。这导致 $W_N$ 对瞬子模空间 $\mathcal{M}_{\mathrm{U}(N)}$ 等变上同调的一个作用；它从 4d 规范理论一侧解释了 Toda CFT 的 $W_N$ 对称性何以出现。
% [Tutor Note: p.102 的"4d 侧解释 $W_N$"是本节结论性的一步，核对路径为：
% 瞬子模空间 $\mathcal{M}_{\mathrm{U}(N)}$ 的等变上同调上有 $\mathrm{SH}^c$（在 $h_1=1$ 的特殊化）的作用，
% 而 $\mathrm{SH}^c$ 在特定参数处等于 $W_N$ 手征代数的泛包络代数，于是 $W_N$ 作用到 4d 的瞬子模空间上；
% 这正是 AGT 中 Toda CFT 的 $W_N$ 对称性在规范理论一侧的来源。]
\end{itemize}

椭圆 Hall 代数、DIM 代数及其众多退化极限，已被用于与 AGT 相关的应用 [36, 281, 282, 385, 642, 648, 651, 929, 944, 970, 978, 980, 1012, 1013, 1017, 1034--1043]。由于这些代数关系密切、对物理应用而言几乎等价，很难理清究竟哪个精确的代数是任一具体物理工作所需要的。选择往往基于代数中哪些生成元在物理上最有意义：从一个代数表述翻译到另一个表述是高度复杂的事情，例如见 [1033, 1044]。

人们也考虑过更多的代数。
\begin{itemize}[leftmargin=1.4em,itemsep=4pt]
  \item \emph{量子环面代数} $\ddot{\mathcal U}_{q_1,q_2,q_3}(\mathfrak{sl}_k)$ 由 [1045--1047] 引入。$\ddot{\mathcal U}_{q_1,q_2,q_3}(\mathfrak{sl}_k)$ 的不同表述之间、它与类似物 $\ddot Y_{h_1,h_2,h_3}(\mathfrak{sl}_k)$ 之间的关系，以及进一步的经典极限，在 [1048--1051] 中得到探讨。我不知道在这一情形下是否已有 DIM 代数或椭圆 Hall 代数的类似物。这一类代数与存在 $\Z_k$ 轨形时的 AGT 对应有关 [1052--1054]，或许还可参见 [37]。
% [Tutor Note: p.102 原文 "I do not know if there is are analogues of the DIM algebra…" 有明显笔误
% （"is are" 应为 "are"），照其义译出，不改内容。]

  \item 正如 $W_N$ 手征代数嵌入 $\mathrm{SH}^c$，$\mathcal{E}_{\sigma,\bar\sigma}$ 的适当特殊化也包含 [1055] $qW$ 代数与箭图 W 代数，后者本身在 [84, 135, 625, 634, 645, 654, 659, 777, 868, 870, 1056--1073] 中得到进一步探讨。

  \item 关于 BPS 代数的一般观点由 COHA 给出 [1066, 1074--1077]。

  \item 相关文献中还出现了 Sklyanin 代数。它是量子群 $U_q(\mathfrak{sl}_2)$ 的一参数形变。
% [Tutor Note: p.102 原文把 Sklyanin 拼作 "Skylanin"（ly 误作 yl），属笔误；此处按正确的 Sklyanin 译出。
% 依据：它是 $U_q(\mathfrak{sl}_2)$ 的一参数形变，即 Sklyanin 于 1982 年引入的 Sklyanin 代数。]
\end{itemize}

基于这些广义对称代数，存在精细化拓扑顶点的一个椭圆版本（椭圆 DIM 代数的交织算子）[927, 1078--1080]，可用于 AGT 对应的 6d 提升；此外还有 Macdonald 精细化拓扑顶点 [656, 660]，以及 4d 时空被轨形化时的类似物 [1053]。

\para{矩阵模型。}
AGT 对应（包括其 $q$-形变版本）可以通过研究矩阵模型来探讨 [954]，因为瞬子配分函数与共形块都是带对数势的 Penner 型矩阵模型积分。综述见 [140, 308]。

在 2d CFT 一侧，矩阵模型表示作为共形块的 Dotsenko--Fateev 自由场表示出现，而后者的可用性取决于每个三穿刺球面中 Liouville/Toda 动量之和是否被适当量子化。
% 说明：末句原文跨页（p.102 页末 "provided the sum of" → p.103 首行），
% 按 STYLE_GUIDE §10「跨页句子开始者写完」由本片段收尾；tr-08 自 p.103 第三行起续写。
